% Notepro template: academic
% Engine: XeLaTeX (Tectonic). A4, wider margins, section numbering, TOC-ready.
% Placeholders: PREAMBLE, TITLE, AUTHOR, DATE, BODY, BIB
\documentclass[11pt,a4paper]{article}
\usepackage[margin=2.5cm]{geometry}
\usepackage{setspace}
\onehalfspacing
\usepackage{amsmath}
\usepackage{amssymb}
% --- CJK (XeLaTeX) ---
\usepackage{ctex}
\usepackage[backend=bibtex,style=numeric]{biblatex}
\addbibresource{references.bib}
\title{相對論筆記 Notes on Relativity}
\author{Ada Lovelace}
\date{\today}
\begin{document}
\maketitle
\begin{abstract}
\end{abstract}
\section{引言 Introduction}

愛因斯坦在 1905 年提出狹義相對論，其質能等價關係為 $E = mc^2$。

考慮以下積分作為一個範例公式：

\[
\int_0^\infty e^{-x^2} \, dx = \frac{\sqrt{\pi}}{2}
\]

這個結果在量子力學中經常出現 \autocite{dirac1930}。

\subsection{結論 Conclusion}

理論物理的優雅之處在於用簡潔的數學描述自然 \autocite[p. 42]{feynman1965}。

\printbibliography
\end{document}
